%% Thesis type: bachelor, master, doctoral
\newcommand*{\ThesisType}{master}

%% Thesis title (exactly as in the formal assignment)
\newcommand*{\ThesisTitle}{Estimation of Quantum Fisher Information by variational quantum circuits on a quantum computer}
%% Plaintext version for PDF metadata, uncomment if needed (defauls to \ThesisTitle)
\newcommand*{\ThesisTitlePlaintext}{TeXtured Manual \TeXturedVERSION}
%% Thesis title (if custom formatting is needed for the title page, defauls to \ThesisTitle)
\newcommand*{\ThesisTitleFront}{
\ThesisTitle\\
{\Huge\color{gray}\TeXturedVERSION}
}
%% Thesis subtitle (optional, uncomment to display a subtitle below the title on the title page)
%%\newcommand*{\ThesisSubtitle}{Sottotitolo}

%% Author of the thesis
\newcommand*{\ThesisAuthor}{Luca Manzi}
%% Plaintext version for PDF metadata, uncomment if needed (defauls to \ThesisAuthor)
\newcommand*{\ThesisAuthorPlaintext}{Luca Manzi}
%% Student ID / matricola number (optional, uncomment to display on title page)
\newcommand*{\ThesisStudentID}{856095}

%% Year when the thesis is submitted
\newcommand*{\YearSubmitted}{2025-2026}
%% Year of the last revision, uncomment if it is different from \YearSubmitted
% \newcommand*{\YearRevision}{2026}

%% University
\newcommand*{\University}{Università degli Studi di Milano Bicocca, Università degli Studi di Milano, Università degli Studi di Pavia}

%% Name of the department or institute, where the work was officially assigned
%% (according to the Organizational Structure of MFF UK in English,
%% or a full name of a department outside MFF)
\newcommand*{\Department}{Dipartimento di Fisica \enquote{Aldo Pontremoli}, Università degli Studi di Milano}

%% Is it a Department (katedra), or an Institute (ústav)?
\newcommand*{\DeptType}{Department}

%% Thesis supervisor: name, surname and titles
\newcommand*{\Supervisor}{Prof. Enrico Prati}
%% Thesis co-supervisor: name, surname and titles (uncomment if applicable)
\newcommand*{\CoSupervisor}{Dr. Sebastiano Corli}

%% Supervisor's department/institute (again according to Organizational structure of MFF)
\newcommand*{\SupervisorsDepartment}{Dipartimento di Fisica \enquote{Aldo Pontremoli}, Università degli Studi di Milano}

%% Study programme and specialization
\newcommand*{\StudyProgramme}{Artificial Intelligence for Science and Technology, \\ Complex systems and Quantum Technologies}

%% Abstract (recommended length around 80-200 words; this is not a copy of your thesis assignment!)
\newcommand*{\Abstract}{%
    Quantum Fisher information (QFI) sets the ultimate sensitivity of a parameterized quantum state, but practical use requires scalable calculations and a model of the complete process.  This Thesis addresses both needs.  It implements the Variational Quantum State Eigensolver within a variational QFI pipeline to estimate bounds without full density-matrix diagonalization.
    A Transverse-Field Ising Model is studied as a quantum magnetometer.  Restricted access reduces the available QFI, and fixed depolarization lowers it across the tested field range.  For $N=6$ and accessible size $n=4$, a depolarizing probability $P=0.3$ reduces the peak local QFI by $36.4\%$.  Finite-temperature Gibbs preparation can instead outperform the reduced ground state away from its optimum.  At $\beta=5$, it provides more local QFI at $40$ of $60$ sampled fields.  The best preparation therefore depends on the expected operating regime.
    The Thesis also studies exoplanet imaging with spatial-mode demultiplexing (SPADE).  Aggregate counts supplied by ASI Matera do not support experimental truncated-QFI estimation or complete channel reconstruction.  A predictive classical calibration model and a modular quantum-to-classical channel sequence are developed instead.  The calibration model achieves about $94\%$ coverage for nominal $95\%$ grouped holdout intervals.  Under the assumed source model, full ideal SPADE keeps at least $99.899\%$ of the source QFI per surviving photon, while an assumed binary cross-talk model can reduce the retained CFI to $0.13\%$--$57.0\%$.  These results motivate a future adaptive VQ-SPADE experiment.
}

%% Subject (short description for PDF metadata)
\newcommand*{\Subject}{%
    Quantum Fisher Information Thesis
}

%% Keywords (about 3-7)
\newcommand*{\Keywords}{%
    Quantum Information, Quantum Fisher Information, QFI, Variational Quantum State Eigensolver, VQSE, Spatial-Mode Demultiplexing, SPADE,
}
%% Plaintext version for PDF metadata, uncomment if needed (defauls to \Keywords)
% \newcommand*{\KeywordsPlaintext}{%
%    Quantum Fisher Information, QFI, Variational Quantum State Eigensolver, VQSE, Spatial-Mode Demultiplexing, SPADE
% }
